% Pillar-4 (length bias) results opener
\section{Length Bias and Generalization: Results}
\label{sec:p4-results}
We report three findings and then scope them. First, in our near-ceiling
Qwen3-8B-Instruct evaluation, GRPO post-training does not deliver a significant
held-out gain: the mean improvement over the pre-RL checkpoint on held-out data is
small ($+0.013$) and not statistically significant (paired $p=0.256$, $n{=}5$). We
stress that this is a property of the near-ceiling regime, not of RL in general---in
the lower-accuracy Qwen2.5-1.5B-Instruct GSM8K-CoT experiment analyzed throughout
the length sections, both GRPO and Dr.\,GRPO \emph{do} produce significant
pre$\to$post held-out gains ($+6.2$pp, McNemar exact $p=1.1\times10^{-4}$;
$+5.0$pp, $p=7.6\times10^{-4}$). Second, GRPO
and Dr.\,GRPO produce held-out gains that are indistinguishable at our seed
budget---removing the length-normalization term does not, here, change the
outcome. Third, under the 200-token GSM8K generation cap, neither algorithm
inflates length: pooled mean completion length drifts \emph{downward}, from
$194.4$ to $183.5$ tokens, and the verbosity-trap signature (a rising
coupling between length and reward as training proceeds) is absent.

Taken together these say that, in this regime, there is no length pathology for
Dr.\,GRPO to fix. The subsections develop the mechanism behind the null---the
weak, often negative length--reward coupling in short arithmetic
chains---and show through lag and reward-shape analyses that length dynamics are
only weakly identifiable here, so the defensible claim is a correction of a formal
gradient bias rather than universal held-out superiority. We then turn the null into
a positive research program: the external frontier-model cross-examination in the
final subsection specifies the length-confounded, sparse-reward regime in which a
Dr.\,GRPO advantage \emph{should} appear (high length dispersion, a train
length-spurious / test length-anti-spurious split), a causal length-mediation
protocol that decomposes held-out success into direct and verbosity-mediated
effects, and a length-adversarial truncation test executable on the checkpoints we
already release.
